% 用户扫描《大学物理简明教程》第5版第9章“变化的电磁场”
% 选择题9.1.1—9.1.13，印刷/PDF第236—237页。
% 题意研究性转述，推导独立；第9.1.11四选项按严格定义均有疑点。

\subsection{选择题：电磁感应、互感及位移电流}

\sourceproblem{9.1.1}{236}{236}
\problemtext 一圆形金属线圈置于均匀且恒定的磁场中，
问下列运动哪些产生感应电流：平移、绕与磁场平行或垂直的直径转动等。
\workings
闭合回路中的感应电动势满足
\[
\mathcal E=-\frac{d}{dt}(\bm B\cdot\bm S)
=-BS\frac d{dt}\cos\theta
\]
（$B,S$恒定）。
纯平移、绕与$\bm B$平行的直径转动都不改变磁通量。
若绕与$\bm B$垂直的直径转动，
面法线与磁场夹角随时间变化，一般有$\mathcal E\ne0$。
故选\textbf{B}。
\crosscheck 不应因“闭合导体有运动”就断定有感应电流：
关键是闭合环路的总磁通量是否随时间改变。

\sourceproblem{9.1.2}{236}{236}
\problemtext 图9.1.2中一根长度$l$的直导线在均匀磁场$\bm B$中以
速度$\bm v$移动，$\bm v$、导线与$\bm B$位于同一纸面，
分别与水平成一定夹角。问导线端点间的动生电动势。
\workings
对刚性平移导体段，
\[
\mathcal E_{ab}=\int_a^b(\bm v\times\bm B)\cdot d\bm l.
\]
$\bm v$和$\bm B$都位于纸面，
故$\bm v\times\bm B$垂直纸面，而导线元$d\bm l$在纸面内，
其点积处处为零。因此
\[
\boxed{\mathcal E_{ab}=0},
\]
原书选\textbf{D}。
\crosscheck 不能把单个夹角$\alpha$生搬进$Blv\sin\alpha$；
该特殊公式还要求明确磁场、导线和速度的三维取向。

\sourceproblem{9.1.3}{236}{236}
\problemtext 关于涡旋电场的四项描述中，
判断“不正确”的一项：涡旋场能对电荷施力、来自变化的磁场、
由静电荷激发、其场线可闭合。
\workings
法拉第--麦克斯韦方程
\[
\nabla\times\bm E_{\mathrm{ind}}
=-\frac{\partial\bm B}{\partial t}
\]
说明磁场的时间变化激发具有非零环流的感应电场，
它对电荷仍有力$q\bm E_{\mathrm{ind}}$，
感应场线在典型对称情形下形成闭合环。
静止电荷激发的库仑静电场是保守场，并非此处所说的涡旋感应场。
故错误项是\textbf{C}。
\crosscheck 对于动态的一般电场可分解静电样部分与非保守部分；
不能把“涡旋电场”误认为“变化电场”本身。

\sourceproblem{9.1.4}{236}{236}
\problemtext 静电场、稳恒磁场、感生电场、随时间变化磁场四类矢量场中，
哪一类是势函数描述下的保守场？
\workings
静电场满足$\nabla\times\bm E_{\mathrm{static}}=\bm0$，
在无奇异、单连通区域可表示为
$\bm E_{\mathrm{static}}=-\nabla\varphi$，
使任意闭路$\oint\bm E_{\mathrm{static}}\cdot d\bm l=0$。
感生电场一般环流非零；磁场本身不按保守力场定义。
故选\textbf{A}。
\crosscheck 此处“保守场”指电场作为单位电荷受力场，
而非仅凭磁场是否随时间改变来判别。

\sourceproblem{9.1.5}{236}{236}
\problemtext 一圆形铜线环和形状完全相同但材料不同的不锈钢线环静置，
穿过两者的磁通量变化率相同，忽略自感，比较感应电动势大小。
\workings
两环均满足法拉第定律
$\mathcal E=-d\Phi_B/dt$，与材料电阻率无关，
因此
\[
\boxed{|\mathcal E_{\mathrm{Cu}}|
=|\mathcal E_{\mathrm{steel}}|},
\]
原书选\textbf{D}。
\crosscheck 材料电阻不同会使感应电流$I=\mathcal E/R$不同，
但题问的是电动势而不是电流。

\sourceproblem{9.1.6}{236}{236}
\problemtext 一匀强磁场中转轴与$\bm B$垂直的线圈匀速转动，
转速增为原来的2倍、其他条件不变。比较最大感应电动势、
电流和交流频率。
\workings
在单匝或固定匝数$N$的转动线圈中，
$\Phi=NBA\cos(\omega t+\phi)$，
\[
\mathcal E=NBA\omega\sin(\omega t+\phi).
\]
最大感应电动势$\mathcal E_{\max}=NBA\omega$，
故转速加倍时最大值加倍，原题选\textbf{A}。
\crosscheck 频率$f=\omega/(2\pi)$同样仅加倍而非四倍；
电流峰值还受电路电阻、电感和频率依赖的阻抗影响，
不能无条件断言其增加四倍。

\sourceproblem{9.1.7}{236}{236}
\problemtext 长$L$的金属杆绕端点$O$以角速度$\omega$在垂直恒定磁场
$\bm B$的平面内匀速转动，求两端动生电动势。
\workings
距轴$r$的杆微元切向速率$v=\omega r$，
局部磁力单位电荷沿杆方向，故
\[
|\mathcal E|=\int_0^L B\omega r\,dr
=\boxed{\frac12B\omega L^2}.
\]
原题选\textbf{E}。
\crosscheck 动生电动势与时间无关；
杆角度$\omega t+\theta$只决定杆在平面中的瞬时朝向，
并不使匀强垂直磁场下的电动势额外乘余弦。

\sourceproblem{9.1.8}{237}{237}
\problemtext 两回路互感分别记为$M_{21}$与$M_{12}$，
且$|di_1/dt|<|di_2/dt|$。比较互感系数及线圈2中由线圈1产生的
电动势$\varepsilon_{21}$与线圈1中由线圈2产生的电动势$\varepsilon_{12}$的大小。
\workings
对线性静止互易电路，互感系数满足
\[
\boxed{M_{12}=M_{21}=M}.
\]
而互感电动势的大小分别为
$|\varepsilon_{21}|=M|di_1/dt|$、
$|\varepsilon_{12}|=M|di_2/dt|$，
故
\[
\boxed{|\varepsilon_{21}|<|\varepsilon_{12}|},
\]
原书选\textbf{D}。
\crosscheck 电动势比较的是绝对值；
有向瞬时电动势符号还取决于线圈绕向、端点参考方向与电流变化符号。

\sourceproblem{9.1.9}{237}{237}
\problemtext 磁能公式$W_m=\tfrac12LI^2$的适用范围，
是否仅限无限长螺线管、单匝线圈、密绕螺线环，或适用于自感恒定的任意线性回路？
\workings
线性磁场能量积累时
$dW_m=I\,d(LI)=LI\,dI$，
在$L$不随电流变化时积分给
\[
\boxed{W_m=\int_0^I L i\,di=\frac12LI^2}.
\]
故适用于自感系数$L$固定的任何线性几何回路，
原题选\textbf{D}。
\crosscheck 对铁磁非线性材料，$L$随$I$可能变化，
一般需积分$W=\int I\,d\Psi$，不能直接写$\tfrac12LI^2$。

\sourceproblem{9.1.10}{237}{237}
\problemtext 真空中一根无限长直导线载流$I$，
问距导线$a$处的磁场能量密度。
\workings
导线场强
$B(a)=\mu_0I/(2\pi a)$，
真空磁能密度$w_m=B^2/(2\mu_0)$，
因此
\[
\boxed{w_m=\frac1{2\mu_0}
\left(\frac{\mu_0I}{2\pi a}\right)^2
=\frac{\mu_0I^2}{8\pi^2a^2}}.
\]
原书选\textbf{B}。
\crosscheck 量纲$[\mathrm{J/m^3}]$；
不要把场能密度错写为$B^2/(2\varepsilon_0)$。

\sourceproblem{9.1.11}{237}{237}
\problemtext 判断位移电流的定义、起因、热效应和磁效应的四种断言：
(A)位移电流就是变化的电场；(B)由变化的磁场产生；
(C)热效应服从焦耳定律；(D)磁效应不服从安培环路定理。
\workings
定义为位移电流密度
\[
\boxed{\bm J_D=\frac{\partial\bm D}{\partial t}}
\]
（真空为$\varepsilon_0\partial\bm E/\partial t$），
其\emph{量纲和物理意义}不是“变化的电场”本身；
时间变化的电场还需乘以$\varepsilon_0$并对时间求导。
变化的磁场按法拉第方程激发涡旋电场，
但不是位移电流\emph{定义}中的必要起因。
理想真空位移电流没有自由载流子欧姆耗散，
不遵循“$I^2R$”的焦耳定律；
它却要计入推广后的安培--麦克斯韦环路定理
\[
\oint_L\bm H\cdot d\bm l
=I_{\mathrm{cond}}+\frac{d}{dt}\int_S\bm D\cdot d\bm S.
\]
所以按四个选项的\emph{严格字面定义}，
\textbf{均不能无条件成立，原题缺少准确选项}。
\crosscheck \textbf{原题待二校：}若教材将(A)宽泛简称为
“位移电流来自变化的电场”，可理解其想选(A)，
但其字面没有量纲一致性，不应把不严谨断言冒充严格定义。

\sourceproblem{9.1.12}{237}{237}
\problemtext 一个充好电的平行板电容器，
左极板$A$正、右极板$B$负，外电路电阻及开关如图9.1.12。
闭合开关后，板间位移电流沿图示$x$轴的哪个方向？
\workings
闭合开关后电容器放电，正负极板电量绝对值逐渐减小。
电场$\bm E$初始从$A$指向$B$，即$+x$方向，
其大小递减，故
$\partial\bm D/\partial t$沿$-x$方向。
因此板间位移电流密度满足
\[
\boxed{\bm J_D=\frac{\partial\bm D}{\partial t},\qquad J_{D,x}<0}.
\]
方向为\textbf{$x$轴负向}，对应原题选\textbf{B}。
\crosscheck 充电时位移电流与电场增长同向；
本题是\emph{放电}，容易误将充电时的方向套用到放电。

\sourceproblem{9.1.13}{237}{237}
\problemtext 平行板电容器\emph{充电}时，
画两条闭合环路$L_1,L_2$，
分别穿过板间的位移电流和引线中的传导电流，
比较$\oint_{L_1}\bm H\cdot d\bm l$、
$\oint_{L_2}\bm H\cdot d\bm l$。
\workings
使用安培--麦克斯韦定律
\[
\oint_L\bm H\cdot d\bm l
=I_{\mathrm{cond,encl}}+I_{\mathrm{disp,encl}}.
\]
同一传输回路充电时，两种电流的\emph{有向数值}大小相等，
$I_D=\varepsilon_0\,d\Phi_E/dt=I$，
只要$L_1,L_2$按同一右手定向约定，
\[
\boxed{\oint_{L_1}\bm H\cdot d\bm l
=\oint_{L_2}\bm H\cdot d\bm l=I}.
\]
原题选\textbf{B}。
\crosscheck 不能只对一条曲面写传导电流而把另一条曲面的位移电流漏掉；
这正是推广安培环路定理的原因。
